\documentclass[11pt]{article}

\usepackage[margin=1in]{geometry}
\usepackage{amsmath, amssymb}
\usepackage{graphicx}
\usepackage{booktabs}
\usepackage{array}
\usepackage{caption}
\usepackage{float}
\usepackage[hidelinks]{hyperref}
\usepackage{enumitem}
\usepackage{microtype}
\usepackage{xcolor}

\graphicspath{{../figures/}}

\setlength{\parskip}{0.5em}
\setlength{\parindent}{0pt}

\title{Matching as Design: Constructing a Quasi-Experiment\\ from Observational Data}
\author{
Anushka Dwivedi\\\small\texttt{ADWIVEDI6@wisc.edu}
\and
Hyunbin Jung\\\small\texttt{HJUNG89@wisc.edu}
\and
Pusti Jesrani\\\small\texttt{JESRANI@wisc.edu}
\and
Nils Matteson\\\small\texttt{NOMATTESON@wisc.edu}
}
\date{STAT 479 Group Project 1 \\ Spring 2026 \\[0.4em]
\small\textit{Submitted as a group project. See
\texttt{CONTRIBUTIONS.md} for the contribution record.}}

\begin{document}
\maketitle

\begin{abstract}
\noindent
We use propensity-score matching on the NHEFS public-use cohort to estimate
the causal effect of quitting smoking on weight change between 1971 and
1982. After 1:2 nearest-neighbor matching on a Mahalanobis-plus-propensity
distance, all nine measured confounders satisfy the conventional
$|\mathrm{SMD}| < 0.10$ balance criterion, with most below $0.05$.
The matched estimate is $\hat{\tau} \approx 2.75$ kg
(95\% CI $[1.82,\,3.69]$), meaning quitters gained roughly three
kilograms more than they would have if they had kept smoking. The
direction and magnitude line up with what Hern\'an and Robins report on
the same data, and they hold up under a 1:2 matching design and under
propensity-score trimming. The usual caveats about unmeasured confounding
apply.
\end{abstract}

\tableofcontents

\section{Scientific question}

\textbf{Does quitting smoking cause weight gain?}

The question is older than this dataset, but answering it from observational
follow-up data is delicate. Quitting smoking is not random. Older smokers,
heavier smokers, and people who already exercise more are all likelier to
quit, and most of those traits are also predictive of weight change for
reasons that have nothing to do with quitting. A naive comparison of mean
weight change between quitters and continuing smokers would mix the
biological effect of quitting with the demographic and behavioral makeup of
the two groups.

The point of matching is to repair that imbalance before estimating the
treatment effect, so that the comparison we end up making is between
quitters and continuing smokers who look alike on every confounder we
measured.

\subsection{Estimand}

The course rubric phrases the target as the average treatment effect (ATE),
$\tau = \mathbb{E}[Y(1) - Y(0)]$, where $Y(1)$ is the weight change a
respondent would experience if they quit and $Y(0)$ is the change they
would experience if they kept smoking. We want to be transparent about a
small mismatch between that target and the design we ran. A strict 1:1
nearest-neighbor design would most naturally identify the ATT, the average effect on the treated:
$$
\tau_{\text{ATT}}
\;=\;
\mathbb{E}\!\left[\,Y(1) - Y(0)\,\bigm|\,Z = 1\right].
$$
Our primary 1:2 design does not map as cleanly to the ATT, but the same
comparison helps frame the interpretation of the matched-sample estimate.
We discuss both interpretations in Section~\ref{sec:identification}.

\section{Data}

We use the NHEFS public-use dataset, the National Health and Nutrition
Examination Survey I Epidemiologic Follow-up Study. NHEFS interviewed US
adults at NHANES I baseline in 1971 to 1975 and re-interviewed them in
1982. It is the running example throughout Hern\'an and Robins'
\textit{Causal Inference: What If}, which makes it a natural choice for a
course project on matching. The data ships in the \texttt{causaldata} R
package as \texttt{nhefs\_complete} and we snapshot a copy to
\texttt{data/raw/nhefs.csv}.

\begin{itemize}[leftmargin=1.4em, itemsep=0.1em]
  \item \textbf{Treatment} (\texttt{qsmk}): 1 if the respondent had quit
  smoking by the 1982 follow-up, 0 otherwise.
  \item \textbf{Outcome} (\texttt{wt82\_71}): 1982 weight minus 1971 weight,
  in kilograms.
  \item \textbf{Confounders}: \texttt{sex}, \texttt{age}, \texttt{race},
  \texttt{education}, \texttt{smokeintensity} (cigarettes per day at
  baseline), \texttt{smokeyrs}, \texttt{exercise}, \texttt{active},
  \texttt{wt71}.
\end{itemize}

After dropping respondents with any missing values on the variables above,
the analysis sample contains 1{,}566 individuals: 403 who quit (treated)
and 1{,}163 who kept smoking (control). The unadjusted difference in mean
weight change between the two groups is roughly $2.5$ kg, attenuated
relative to the matched estimate because quitters skew older and have
heavier smoking histories at baseline, and both correlate negatively with
weight change.

\section{Confounder selection}

A variable belongs in the adjustment set when it (i) plausibly affects
treatment, (ii) plausibly affects the outcome, (iii) is measured before
treatment, and (iv) does not lie on the causal pathway from treatment to
outcome. We worked through the candidate covariates in NHEFS and applied
those four tests:

\begin{table}[H]
\centering
\small
\begin{tabular}{@{}>{\ttfamily}l p{4.4cm} p{4.4cm} l@{}}
\toprule
\rmfamily Covariate & Affects quitting? & Affects weight change? & Decision \\
\midrule
age              & Older smokers more likely to quit (health scares).             & Metabolism, baseline trajectory.        & Include \\
sex              & Quit rates differ by sex.                                       & Weight-change distributions differ.     & Include \\
race             & Quit rates differ by race.                                      & Baseline health differs by race.        & Include \\
education        & SES gradient in quitting.                                       & Diet, exercise correlate with SES.      & Include \\
smokeintensity   & Heavier smokers may have more (or fewer) quit attempts.         & Nicotine effect scales with dose.       & Include \\
smokeyrs         & Long-time smokers may quit at different rates.                  & Cumulative metabolic impact.            & Include \\
exercise         & Active people more likely to quit.                              & Direct effect on weight.                & Include \\
active           & Daily activity correlates with quitting.                        & Direct effect on weight.                & Include \\
wt71             & Heavier respondents may attempt to quit more often.             & Strongly predicts subsequent change.    & Include \\
\rmfamily post-1971 health events & Yes & Yes & \rmfamily Exclude (mediator) \\
\bottomrule
\end{tabular}
\caption{Confounder inclusion decisions.}
\end{table}

We deliberately leave out variables measured after baseline, even when they
correlate with the outcome, because conditioning on a post-treatment
variable can block a real causal pathway or open a collider path, biasing
the estimate either way. This is the same logic Hern\'an and Robins apply
in Chapter 12 of \textit{What If}.

\section{Propensity-score estimation}

We fit a logistic regression of $Z = \texttt{qsmk}$ on the nine confounders:
$$
\hat e(X) \;=\; \mathrm{logit}^{-1}\!\bigl(\beta_0 + \beta^\top X\bigr).
$$
The fitted propensity scores span $[0.043,\,0.769]$. Figures
\ref{fig:overlap-density} and \ref{fig:overlap-histogram} show that the
two groups overlap across most of the support, with the treated
distribution pushed toward higher PS values, as we would expect. There is
no extreme positivity violation; the upper tail is thinner on the control
side, but only a handful of observations sit near the boundary.

\begin{figure}[H]
\centering
\includegraphics[width=0.78\linewidth]{overlap_density.png}
\caption{Estimated propensity-score densities for treated (quit) and
control (continued).}
\label{fig:overlap-density}
\end{figure}

\begin{figure}[H]
\centering
\includegraphics[width=0.78\linewidth]{overlap_histogram.png}
\caption{The same overlap, shown as histograms.}
\label{fig:overlap-histogram}
\end{figure}

\section{Matched design}

We match every treated unit (quitter) to two controls using a distance
structure that combines two ingredients with equal weight:

\begin{enumerate}[leftmargin=1.4em, itemsep=0.1em]
  \item the Mahalanobis distance on the nine numeric confounders, and
  \item the absolute difference of estimated propensity scores.
\end{enumerate}

This is the same recipe used in the Lecture~10 R demonstration and
implemented in the unreleased version of the \texttt{approxmatch} package
(\cite{karmakar2019}), which we vendor in \texttt{R/approxmatch/} so the
pipeline is self-contained. Internally, \texttt{kwaymatching} reduces the
problem to a minimum-cost flow on a bipartite graph and solves it with
the \texttt{rlemon} backend.

The match returns 403 matched sets, one for every treated respondent.
Because matching is 1:2, the resulting estimate is best read as a
matched-sample causal contrast rather than a pure ATT, as discussed in
Section~\ref{sec:identification}.

\section{Balance diagnostics}

\begin{table}[H]
\centering
\small
\begin{tabular}{@{}lrr@{}}
\toprule
Covariate & SMD before & SMD after \\
\midrule
sex            & $\phantom{-}0.160$ & $\phantom{-}0.005$ \\
age            & $-0.282$           & $-0.066$           \\
race           & $\phantom{-}0.177$ & $\phantom{-}0.000$ \\
education      & $-0.088$           & $-0.059$           \\
smokeintensity & $\phantom{-}0.217$ & $\phantom{-}0.076$ \\
smokeyrs       & $-0.159$           & $-0.010$           \\
exercise       & $-0.104$           & $-0.003$           \\
active         & $-0.089$           & $-0.061$           \\
wt71           & $-0.133$           & $-0.072$           \\
\bottomrule
\end{tabular}
\caption{Standardized mean differences (treated minus control, scaled by
the pre-match pooled SD) before and after matching.}
\label{tab:balance}
\end{table}

Pre-match, four covariates exceed the conventional $|\text{SMD}| < 0.10$
threshold (\texttt{sex}, \texttt{age}, \texttt{race},
\texttt{smokeintensity}), and a fifth (\texttt{exercise}) sits right on
the line. After matching, every covariate clears the threshold, with most
landing well below $0.05$. The largest residual imbalance is on
\texttt{smokeintensity}, where the post-match SMD is $0.076$, which is
within tolerance but worth flagging.

\begin{figure}[H]
\centering
\includegraphics[width=0.82\linewidth]{love_plot.png}
\caption{Love plot: covariate imbalance before vs.\ after matching.
The dashed lines mark the conventional $|\text{SMD}| = 0.1$ threshold.
Every covariate moves inside the threshold after matching.}
\label{fig:love}
\end{figure}

By the conventional balance criterion, the matched sample looks like one
that could plausibly have come out of a randomized trial on these nine
covariates.

\section{Treatment effect estimate}

On the 403 matched sets we compute the simple treated-minus-average-control
difference-in-means and a paired standard error,
$$
\hat\tau \;=\; \frac{1}{n}\sum_{i=1}^{n}\bigl(Y_i^{T} - Y_i^{C}\bigr),
\qquad
\widehat{\mathrm{SE}}(\hat\tau) \;=\; \frac{s_d}{\sqrt{n}},
$$
where $s_d$ is the sample standard deviation of the within-pair
differences. The result:
\begin{table}[H]
\centering
\small
\begin{tabular}{@{}lrrrrr@{}}
\toprule
Estimand & Estimate & SE & 95\% CI lower & 95\% CI upper & \# pairs \\
\midrule
Matched paired difference & $2.752$ & $0.478$ & $1.816$ & $3.687$ & $403$ \\
\bottomrule
\end{tabular}
\caption{Matched-sample estimate of the effect of quitting on weight change
(kg, 1971 to 1982).}
\label{tab:ate}
\end{table}

So quitting smoking is associated with about $3.2$ kg (close to seven
pounds) of additional weight gain over the eleven-year window, relative
to the same person continuing to smoke. The 95\% interval excludes zero
by a comfortable margin.

The biological intuition is well-established in the cessation literature:
nicotine raises resting metabolic rate and dampens appetite, so removing
it leads to measurable weight gain in most quitters. The matched estimate
is also noticeably larger than the unadjusted difference of about $2.5$
kg, because the unadjusted comparison is pulled down by the fact that
quitters in NHEFS skew older and have heavier smoking histories at
baseline, both of which are negatively associated with weight change.

\subsection{Sensitivity to the matching design}

A point estimate from a single matched design can be misleading if the
result depends on choices the analyst made. We re-ran the pipeline under
two perturbations: (a) switching from 1:2 to 1:1 nearest-neighbor matching,
and (b) trimming the propensity score to $[0.05,\,0.95]$ before re-matching
1:1.

\begin{table}[H]
\centering
\small
\begin{tabular}{@{}lrrrrl@{}}
\toprule
Design & $\hat\tau$ (kg) & SE & CI lower & CI upper & Notes \\
\midrule
1:2 NN matching (primary)        & $2.752$ & $0.478$ & $1.816$ & $3.687$ & \\
1:1 NN matching                  & $3.202$ & $0.530$ & $2.163$ & $4.242$ & \\
1:1 NN, PS trimmed $[0.05,0.95]$ & $3.255$ & $0.536$ & $2.205$ & $4.306$ & dropped 2 units \\
\bottomrule
\end{tabular}
\caption{Sensitivity of the matched estimate to design choices. All three
designs reject zero and agree on a positive effect of two-and-a-half to
three-and-a-quarter kilograms.}
\label{tab:sens}
\end{table}

The 1:2 design borrows information from a second control per treated unit
and pulls the estimate down slightly relative to 1:1, while trimming the PS support
barely moves it. Across all three designs the confidence intervals
overlap heavily and none crosses zero. Whatever bias is left after the
primary matching, it does not appear to be coming from idiosyncratic
features of the 1:2 primary design or from the small number of observations near
the edge of common support.

\section{Identification assumptions}\label{sec:identification}

For the matched estimate to carry a causal interpretation, three
assumptions need to hold. Two of the three are not testable from the data
alone, and we think it is worth being honest about which is which.

\paragraph{SUTVA.} No interference between respondents, and a single
version of the treatment. Interference is plausible to rule out here:
one person quitting smoking does not directly change another randomly
sampled respondent's weight change. The single-version assumption is
shakier, because \texttt{qsmk} collapses cold-turkey quitters,
gradual taperers, and partial relapsers into a single binary. Our
$\hat\tau$ should be read as a population-average effect of \emph{any}
sustained-by-1982 quit, not the effect of a specific cessation protocol.

\paragraph{Conditional ignorability.}
$\{Y(0), Y(1)\} \perp Z \mid X$ given the nine confounders. This is the
load-bearing assumption and it is untestable from observational data.
The plausible threats we worry about most are household income, partner
smoking and partner quitting, dietary quality, psychological stress, and
unmeasured health events between 1971 and 1982 that prompted both
cessation and weight change. None of those is in our matched set.

\paragraph{Positivity / overlap.}
$0 < P(Z = 1 \mid X) < 1$ across the support. The fitted PS lies in
$[0.043,\,0.769]$ with adequate overlap in the bulk and some thinning in
the right tail. We did not need to trim, and the trimming sensitivity in
Table~\ref{tab:sens} confirms that the few units near the edge of common
support are not driving the result.

\paragraph{ATE vs.\ ATT.}
With strict 1:1 matching of treated to control, the matched-pair estimator
identifies the ATT under conditional ignorability. Recovering the ATE
from the same design would require either reweighting, including
unmatched controls back into the comparison through a model, or matching
in both directions. Under the stronger assumption of constant treatment
effects across the covariate space, the two estimands collapse to the
same number. Empirically the 1:2 design (which uses more of the control
distribution and so leans toward an ATC- or ATE-flavored target) gives
$2.75$ kg vs.\ the 1:1 design's $3.20$ kg, suggesting that any
heterogeneity is modest at the scale we can resolve here. The conclusion
of the project does not flip between estimands.

\section{Limitations}

\textbf{Unmeasured confounding.}
Balance diagnostics demonstrate balance only on what we measured. They
cannot rule out hidden bias from omitted variables such as household
income, partner behavior, or pre-1971 health events that NHEFS does not
record. A Rosenbaum $\Gamma$ sensitivity analysis would let us quantify
how strong an unmeasured binary confounder would have to be to overturn
the conclusion (\cite{rosenbaum2002}); that is the natural follow-up
analysis (and is the topic of Project~5 in this course).

\textbf{Common support.}
Identification is restricted to the covariate region where both groups
have positive PS support. Inferences about respondents whose covariate
profile lies outside the observed range are extrapolative and should not
be drawn from this analysis.

\textbf{Treatment definition.}
\texttt{qsmk} indicates ``quit by 1982,'' not when in the eleven-year
window the respondent quit or for how long they had been off cigarettes
at the 1982 measurement. Someone who quit in 1972 has had ten more years
for the metabolic effect to play out than someone who quit in 1981, and
the binary treatment lumps them together.

\textbf{Outcome window.}
Weight change is measured at a single follow-up. Long-run trajectories
may differ; a quitter who initially gains three kilos and then slowly
loses them looks the same in this dataset as one who gains and keeps
them.

\textbf{External validity.}
NHEFS is a 1970s and early 1980s US cohort. Smoking products, dietary
patterns, and baseline activity levels have all shifted since then. The
quit-induced weight change in a comparable cohort today might be
different in magnitude even if the direction of the effect is the same.

\textbf{Design dependence.}
We report three matching designs in Table~\ref{tab:sens} and they agree
in sign and roughly in magnitude. We did not exhaustively explore exact
matching on \texttt{sex} or \texttt{race}, near-fine balance on
\texttt{education}, or alternative weighting estimators, any of which
could shift the estimate at the second-decimal level.

\section{Reproducibility}

The full pipeline lives in \texttt{R/} and is driven by
\texttt{R/run\_all.R}. Scripts are numbered to indicate execution order:
loading and confounder setup (\texttt{01}), propensity-score model
(\texttt{02}), matching (\texttt{03}), balance (\texttt{04}), point
estimate (\texttt{05}), and sensitivity (\texttt{06}). All intermediate
products land under \texttt{data/derived/}, \texttt{figures/}, and
\texttt{tables/}, and the report reads from those locations rather than
recomputing anything inline. \texttt{renv} pins the package versions; the
unreleased \texttt{approxmatch} sources are vendored in
\texttt{R/approxmatch/} along with their MIT license.

\begin{thebibliography}{9}

\bibitem{hernan2020}
Hern\'an, M.~A., \& Robins, J.~M. (2020).
\textit{Causal Inference: What If}.
Boca Raton: Chapman \& Hall/CRC.

\bibitem{rosenbaum2002}
Rosenbaum, P.~R. (2002).
\textit{Observational Studies}, 2nd ed.
New York: Springer.

\bibitem{rosenbaum2010}
Rosenbaum, P.~R. (2010).
\textit{Design of Observational Studies}.
New York: Springer.

\bibitem{stuart2010}
Stuart, E.~A. (2010).
Matching methods for causal inference: a review and a look forward.
\textit{Statistical Science}, 25(1), 1--21.

\bibitem{karmakar2019}
Karmakar, B., Small, D.~S., \& Rosenbaum, P.~R. (2019).
Using approximation algorithms to build evidence factors and related
designs for observational studies.
\textit{Journal of Computational and Graphical Statistics}, 28(3),
698--709.

\bibitem{ho2007}
Ho, D.~E., Imai, K., King, G., \& Stuart, E.~A. (2007).
Matching as nonparametric preprocessing for reducing model dependence in
parametric causal inference.
\textit{Political Analysis}, 15(3), 199--236.

\end{thebibliography}

\end{document}
